\section{结论}\label{sec:conclusion}
本文研究同步单天线站对未知宽带强弱重叠源的直接定位，以最小二乘估计消去波形系数，将问题化为位置优化。多候选联合AP保留少量粗搜索候选，分别联合调整源位置，再按最终残差选择结果；波形消元块NOMP在每次局部位置试探时重新估计系数，减少固定系数所造成的重复小步。

在2400条正式观测中，五候选联合AP达到99\%的双源20\,m成功率，平均求解时间为0.718\,s；相对单候选恢复96条观测，强源精度保持接近。单候选的共同精度配对显示，联合精修将平均耗时从0.439\,s降至0.396\,s，同时减少目标与梯度计算。共同240条比较中，五候选联合AP、充分AP和宽带SBL均恢复全部双源，而完整停止成本存在明显差异。时间轨迹与典型案例表明，保留候选能够减少错误区域选择，联合更新则减少完成精修的计算。曲率推导为波形消元更新提供了进一步的机制解释。后续可在这一框架下研究未知源数和更一般的传播条件。
